\documentclass{beamer}
\usepackage{ctex}
\usepackage{graphicx,booktabs}
\usepackage{tikz}
\usepackage{comment}
\usepackage{multirow}
\usetikzlibrary{arrows,arrows.spaced,arrows.meta,calc,intersections,through,backgrounds,math,angles,shapes}
\usetheme{Boadilla}
\beamersetaveragebackground{black!6}
\def\mathfamilydefault{\rmdefault}
\setbeamersize{text margin left=1.3cm}\setbeamersize{text margin right=1.2cm}
\newcommand\mytitle[1]{\title{#1}\author{王崇宁}\institute[四中]{郑州四中}\logo{\includegraphics[width=1cm]{cows.jpg}}\maketitle\kaishu\Large}
\newcommand\colorwordsa[1]{\textcolor[rgb]{0,.63,.91}{\heiti #1}}
\newcommand\colorwordsb[1]{\textcolor[rgb]{.75,.53,.21}{\heiti #1}}
\setbeamercolor{mycolor}{fg=black,bg=pink}
\setbeamercolor{myupcol}{fg=white,bg=purple}
\setbeamercolor{mylowcol}{fg=black,bg=pink}
\newcommand\beamercolortextbox[1]{\settowidth\lengtha{#1}\hspace{.5em}\raisebox{-.6ex}{\begin{beamercolorbox}[rounded=true,shadow=true,wd=\lengtha,colsep*=-.4ex]{mycolor}#1\end{beamercolorbox}}}
\newcommand\beamercolorlinebox[1]{\settowidth\lengtha{#1}\centerline{\begin{beamercolorbox}[rounded=true,shadow=true,wd=\lengtha]{mycolor}#1\end{beamercolorbox}}}\newlength\lengtha
\newcommand\threesides{\begin{picture}(14,12)\put(0,0){\qbezier(0,0)(0,8)(7,12)\qbezier(7,12)(14,8)(14,0)\qbezier(14,0)(7,-4)(0,0)}\end{picture}}


\def\displaygray{\setbeamercovered{transparent}}
\newcommand*\circled[1]{\hspace{1pt}\tikz[baseline=(char.base)]\node[very thin,shape=circle,draw,inner sep=.5pt,minimum size=2pt](char){#1};\hspace{1pt}}%



\newcommand\Dd{\displaystyle}\newcommand{\Tt}{\textstyle}\newcommand{\Ss}{\scriptstyle}\newcommand{\Sss}{\scriptscriptstyle}%

\setlength\mathsurround{0.5ex}%文本公式前后留下空白
\makeatletter
\renewcommand\normalsize{\@setfontsize\normalsize\@xpt\@xiipt\abovedisplayskip8\p@\@plus3\p@\@minus5\p@\abovedisplayshortskip\z@\@plus3\p@\belowdisplayshortskip5\p@\@plus3\p@\@minus7\p@\belowdisplayskip\abovedisplayskip\let\@listi\@listI}%行间公式上下
\@addtoreset{equation}{section}
\renewcommand\theequation{\oldstylenums{\thesection}.\oldstylenums{\arabic{equation}}}%公式按章节编号
\makeatother
%\setlength\mathsurround{0.ex}%数学模式与文本模式混排时留出的间距

\newcommand\an[1][a]{\ensuremath{\{#1_n\}}}%sequence
\newcommand\ud{\mathrm{d}}
\newcommand{\ue}{\mathrm{e}}%正体字母
\newcommand\triabc{\ensuremath{\triangle ABC}}%三角形ABC
\newcommand\cnm[2][n]{\ensuremath{\textrm{C}^{#2}_{#1}}}%组合数
%数学题目编辑
\newcommand\lines[1][1.2]{\,\underline{\mbox{\hspace{#1cm}}}\,}% 填空题的横线
\newcommand\brackets[1][2]{\nolinebreak\hfill\mbox{~(\hspace{#1em})}\\}% 选择题的括号
%扩展命令
\newcommand\qqquad{\qquad\quad}
\def\aside#1{\marginpar{\footnotesize #1}}
%自动编号之最简
\newcommand\numa{\refstepcounter{numi}\thenumi}\newcounter{numi}
\newcommand\numb{\refstepcounter{numii}\thenumii}\newcounter{numii}
\newcommand\numc{\refstepcounter{numiii}\thenumiii}\newcounter{numiii}
\newcommand\numaa{\refstepcounter{numai}\thenumai}\newcounter{numai}[numi]

\newcommand{\question}[1][]{\par\vspace{1ex}\noindent\refstepcounter{numberi}\textbf{\thenumberi.}\ensuremath{#1}}\newcounter{numberi}[subsection]% 每道小题自动编号
\newcommand\quson{\\ \hspace*{1em}\refstepcounter{numberii}\thenumberii)~}\newcounter{numberii}[numberi]% 每道小题自动编号
\newcommand\choice[5][4]{\vspace*{-1em}\begin{tasks}(#1)\task$#2$\task$#3$\task$#4$\task$#5$\end{tasks}\vspace*{-1.2em}}%选择题排版之数学模式
\newcommand\choicex[5][4]{\vspace*{-1em}\begin{tasks}(#1)\task#2\task#3\task#4\task#5\end{tasks}\vspace*{-1.2em}}%选择题排版之yysg模式
\newcommand\tbs[1][]{\texttt{\char92#1}}

\newcommand\set[1]{\lbrace\ensuremath{#1}\rbrace}
\newcommand\setx[1]{\{#1\}}

\newcommand\mybf[1]{{\bm#1}}
\def\bfR{\mybf R}
\def\bfN{\mybf N}
\def\bfZ{\mybf Z}
\def\bfQ{\mybf Q}
\def\bfC{\mybf C}
\def\bfZp{\mybf Z^+}

\newcommand\myvec[1]{\bm#1}
\def\veca{\myvec a}
\def\vecb{\myvec b}
\def\vecc{\myvec c}
\def\vecd{\myvec d}
\def\vece{\myvec e}
\def\vecf{\myvec f}
\def\vecs{\myvec e}
\def\vecp{\myvec p}
\def\vecq{\myvec q}
\def\veci{\myvec i}
\def\vecj{\myvec j}
\def\vecei{\myvec{e_1}}
\def\veceii{\myvec{e_2}}
\def\veczero{\myvec 0}
\newcommand\lvec[1]{\overrightarrow{#1}}





\begin{document}
\kaishu
\mytitle{有理数的加法}
\begin{frame}{比一比:看谁读得好}
  \kaishu
  \begin{block}{学习目标：}
    \begin{enumerate}
      \item 经历探索有理数加法法则的过程，体会探索数学问题的乐趣。
      \item 记住有理数的加法法则,能熟练运用有理数的加法法则进行运算.
    \end{enumerate}
  \end{block}
\end{frame}


\begin{frame}{有理数的加法: 探索过程我参与}
  \begin{exampleblock}{同学们配合很好}  
    \begin{enumerate}
      \item +6 (带有负数标记的从前往后站起来6个人)  \pause
      \item -7 (带有正数标记的从前往后站起来7个人)
    \end{enumerate}
  \end{exampleblock}
\end{frame}

\def\hspacea{\hspace*{2in}}

\begin{frame}{有理数的加法: 探索过程我参与}  \pause
  \begin{block}{两个正有理数相加}  \pause
    \hspacea $(+5)+(+7)$  \pause
  \end{block}
  \begin{block}{两个负有理数相加}  \pause
    \hspacea $(-5)+(-7)$  \pause
  \end{block}
  \begin{alertblock}{总结}
     同号两数相加, 取相同的符号, 并把绝对值相加. 
  \end{alertblock}
\end{frame}
\begin{frame}{两个同号的有理数相加}
  \begin{exampleblock}{例2.}
    \begin{enumerate}[<+->]
      \item $98+101$=$\lines$.
      \item $(-4.35)+(-5.45)$=$-(\lines+\lines)$=$-\lines$.
      \item $(-\frac52)+(-2.5)$=$-(\lines+\lines)$=$-\lines$.
    \end{enumerate}
  \end{exampleblock}
\end{frame}

\begin{frame}{有理数的加法: 探索过程我参与}
    \begin{block}{一个正数加一个负数}  \pause
    \hspacea $4+(-3)$
  \end{block}  \pause
  \begin{block}{一个负数加一个正数} \pause
   \hspacea $(-9)+6$
  \end{block}  \pause
  \begin{block}{互为相反数的两个数相加} \pause
    \hspacea $(-3)+3$
  \end{block}  \pause
  \begin{alertblock}{总结}
    异号两数相加,绝对值相等时和为$0$；绝对值不等时, 取绝对值较大的数的符号, 并用较大的绝对值减去较小的绝对值.
  \end{alertblock}
\end{frame}

\begin{frame}{两个异号的有理数相加}  \pause
  \begin{exampleblock}{例1.}
    \begin{enumerate}[<+->]
      \item $(-9)+17$ 的符号是\lines, 绝对值是\lines, 结果是\lines.
      \item $(-\frac83)+2$ 的符号是\lines, 绝对值是\lines, 结果是\lines.
      \item $98+(-101)$=$-(\lines-\lines)$=$-\lines$.
      \item $(-4.35)+5.45$=$+(\lines-\lines)$=$\lines$.
    \end{enumerate}
  \end{exampleblock}
\end{frame}


\begin{frame}{有理数加法法则}
  \kaishu \displaygray
  \begin{block}{}
    \begin{itemize}
      \item<2->  同号两个数相加, 取相同的符号, 并把绝对值相加. %\par
      \item<3-> 异号两个数相加, 绝对相等时和为$0$, 绝对值不等时, 取绝对值较大的数的符号, 并用较大的绝对值减去较小的绝对值. %\par
      \item<4> 一个数同$0$相加, 仍得这个数. 
    \end{itemize}
  \end{block}
\end{frame}

\def\vspacea{\vspace*{-4em}}
\begin{frame}{加法法则的运用}
  \begin{exampleblock}{例3.计算下列各题:}
    (1)$180+(-10)$;\quad   (2)$(-10)+(-1)$;\quad  (3)$5+(-5)$;\quad (4)$0+(-2)$;
  \end{exampleblock}  \pause
  \begin{block}{解:}  \pause
    \begin{columns}  
      \begin{column}{0.5\textwidth}  
        \begin{align*}
          (1)\ &180+(-10)\\  
               &=+(180-10)\\
               &=170;  
        \end{align*}
        \vspacea
        \begin{align*}
          (2)\ &(-10)+(-1)\\ 
               &=-(10+1)\\  
               &=-11;  
        \end{align*}
      \end{column}
      \begin{column}{0.5\textwidth}
        \begin{align*}
          (3)\ &5+(-5)\\  
              &=0\\  
        \end{align*}
        \vspacea
        \begin{align*}
          (4)\ &0+(-2)\\ 
               &=-2\\  
        \end{align*} 
 
      \end{column}
    \end{columns}
  \end{block}
\end{frame}

\begin{frame}{加法法则的运用}
  \begin{block}{练习1.计算下列各题:}
    (1)$-25+(-7)$;\quad   (2)$(-13)+5$;\quad  (3)$(-23)+0$;\quad (4)$45+-(45)$;
  \end{block}
\end{frame}


\begin{frame}
\includegraphics[width=4in]{COWS.jpg}
\end{frame}
\end{document}
